\documentclass[11pt]{article}
\usepackage[utf8]{inputenc}
\usepackage{ngerman}
\usepackage[T1]{fontenc}
\usepackage{geometry}
\geometry{a4paper, margin=2.5cm}
\usepackage{enumitem}
\usepackage{longtable}
\usepackage{booktabs}
\usepackage{amssymb}
\usepackage{tikz}
\usetikzlibrary{shapes,arrows,positioning,fit,backgrounds}

\title{Objektiv-hermeneutische Sequenzanalyse von Text 4}
\subtitle{Gemüsestand, Aachener Markt, 28. Juni 1994}
\author{Paul Koop}
\date{}

\begin{document}

\maketitle

\section*{Protokolleichen}

\begin{tabular}{ll}
\texttt{laanngsaam} & gedehntes Sprechen\\
\texttt{LANGSAM} & lautes Sprechen\\
\texttt{A: ... B: ...} & gleichzeitiges Sprechen\\
\texttt{K} & Käufer\\
\texttt{V} & Verkäufer\\
\texttt{(Kommentare)} & Kommentar des Autors\\
\texttt{. . .} & Pause\\
\texttt{*** Anfang Text ***} & Markierung des Textanfangs\\
\texttt{*** Ende Text ***} & Markierung des Textendes\\
\end{tabular}

\newpage

\section{Sequenzanalyse im Wortlaut}

\subsection*{1. *** Anfang Text4 ***}

\textbf{Paraphrase:} Ein Anfang eines Textes wird markiert.

\textbf{Intention:} Der Leser soll wissen, dass hier ein neuer Text beginnt.

\textbf{Konstitutive Regel:} Ein Text ist nur ein lesbarer Text, wenn er Anfang und Ende hat.

\subsection*{2. : Markt, 11.00 Uhr (Aachen, 28.06.94, Gemüsestand)}

\textbf{Paraphrase:} Der Text bezieht sich auf ein Ereignis auf einem Markt, an einem bestimmten Ort zu einer bestimmten Zeit.

\textbf{Intention:} Der Text erhält eine historische und funktionale Einordnung. Die Ereignisse spielen sich im Rahmen der sozialen Institution Markt ab.

\textbf{Regulative Regel:} Nur zahlungsfähige Kunden erhalten Ware.

\textbf{Konstitutive Regel:} In einer arbeitsteiligen Gesellschaft gibt es ohne Markt und ohne Zahlungsmittel keinen Güteraustausch.

\subsection*{3. (unverständlich)}

\textbf{Paraphrase:} Das Gespräch ist unverständlich. Es läuft aber schon. Der weitere Text bezieht sich auf ein schon laufendes Gespräch.

\textbf{Intention:} Es soll deutlich werden, dass ein laufendes Gespräch verfolgt wird.

\subsection*{4. FK1.1: Hör ens, ich nehm ein paar Champignons mit.}

\textbf{Paraphrase:} Eine Käuferin spricht eine andere Person in der Du-Form an und teilt ihr mit, dass sie einige Champignons mitnimmt.

\textbf{Intentionen (mehrere Lesarten):}
\begin{itemize}
    \item Sie will die Ware einfach nehmen.
    \item Sie will die Ware testen/probieren.
    \item Sie will die Ware kaufen.
\end{itemize}

\subsection*{5. FV1.1: Eh, braune?}

\textbf{Paraphrase:} Die Verkäuferin fragt, ob die Käuferin braune Champignons wünscht.

\textbf{Intention:} Die Verkäuferin will den Bedarf klären.

\textbf{Regulative Regel:} Verkäufer fragen bei Unklarheit nach.

\textbf{Konstitutive Regel:} Ein Verkaufsgespräch erfordert Bedarfsklärung.

\subsection*{6. FK1.2: Ne helle.}

\textbf{Paraphrase:} Die Käuferin sagt, dass sie helle Champignons wünscht.

\textbf{Intention:} Die Käuferin spezifiziert ihren Bedarf.

\subsection*{7. FV1.2: Helle.}

\textbf{Paraphrase:} Die Verkäuferin bestätigt, dass die Käuferin helle Champignons wünscht.

\textbf{Intentionen:}
\begin{itemize}
    \item Die Verkäuferin will der Käuferin helle Champignons einfach geben.
    \item Die Verkäuferin will der Käuferin helle Champignons zum Probieren überlassen.
    \item Die Verkäuferin will der Käuferin helle Champignons verkaufen.
\end{itemize}

\subsection*{8. FK1.3: mhmh.}

\textbf{Paraphrase:} Die Käuferin bestätigt, dass sie helle Champignons wünscht.

\textbf{Intention:} Sie veranlasst die Verkäuferin, ihr helle Champignons zu geben / zu verkaufen.

\subsection*{9. (unverständlich)}

\subsection*{10. FK1.4: Meinen se nich.}

\textbf{Paraphrase:} Die Käuferin hat in die Sie-Form gewechselt und fragt, ob die Verkäuferin helle oder braune Champignons besser findet.

\textbf{Intentionen:}
\begin{itemize}
    \item Sie will die Ware einfach nehmen (wird mit Wechsel in Sie-Form weniger plausibel).
    \item Sie will die Ware testen.
    \item Sie will die Ware kaufen.
\end{itemize}

\subsection*{11. FV1.3: Ja is ejal, se sinn beide frisch.}

\textbf{Paraphrase:} Die Verkäuferin sagt, dass es egal sei, beide Sorten seien frisch.

\textbf{Intention:} Die Verkäuferin will die Entscheidung erleichtern / Bedenken zerstreuen.

\subsection*{12. FK1.5: Oder, wie is et denn mit, mit, eh}

\textbf{Paraphrase:} Die Käuferin ist unsicher geworden. Sie scheint einen neuen Wunsch gebildet zu haben.

\textbf{Regulative Regel:} Ein Kunde darf in Grenzen seinen Kaufwunsch im Gespräch ändern.

\textbf{Konstitutive Regel:} Nur Verkaufsgespräche, die flexibel auf Unsicherheiten reagieren, sind erfolgreich.

\emph{Ab hier werden alle Lesarten, die sich nicht auf einen Kaufwunsch beziehen, aufgegeben.}

\subsection*{13. FV1.4: Die können se länger liejen lassen.}

\textbf{Paraphrase:} Die Verkäuferin liefert einen Grund für die Wahl heller Champignons: Sie lassen sich länger lagern.

\textbf{Intentionen:}
\begin{itemize}
    \item Die Verkäuferin will einen echten Rat geben.
    \item Die Verkäuferin will ihren vorgetäuschten Rat wasserdicht gegen Aufdeckung absichern.
\end{itemize}

\textbf{Regulative Regel:} Verkäuferinnen müssen auf Unsicherheiten eingehen.

\textbf{Konstitutive Regel:} Verkaufsgespräche sind auch Beratungsgespräche.

\subsection*{14. FK1.6: Neh, aber Pfifferlinge.}

\textbf{Paraphrase:} Die Käuferin fragt nach Pfifferlingen.

\textbf{Intention:} Sie will Pfifferlinge kaufen.

\textbf{Regulative Regel:} Ein Käufer muss seinen Kaufwunsch signalisieren.

\textbf{Konstitutive Regel:} Ohne Kaufwunsch kein Kauf.

\subsection*{15. FV1.5: Ah, die sinn super.}

\textbf{Paraphrase:} Die Verkäuferin stellt begeistert fest, dass die Pfifferlinge super sind.

\textbf{Intention:} Sie will das Produkt positiv hervorheben / den Kaufwunsch bestärken.

\subsection*{17. FK1.7: Kann ich die denn in Reissalat tun?}

\textbf{Paraphrase:} Die Käuferin fragt, ob sie die Pfifferlinge in Reissalat mitverarbeiten kann.

\textbf{Intention:} Sie möchte wissen, ob das geht (Verwendungsfrage).

\textbf{Regulative Regel:} Käufer dürfen im Verkaufsgespräch Fragen stellen.

\textbf{Konstitutive Regel:} Verkaufsgespräche sind nur Verkaufsgespräche, wenn Fragen der Käufer zugelassen sind.

\subsection*{19. FK1.8: Brauch ich nich abzukochen oder was?}

\textbf{Paraphrase:} Die Käuferin fragt, ob sie die Pfifferlinge abkochen muss.

\textbf{Intentionen:}
\begin{itemize}
    \item Sie will wissen, ob es gefährlich ist, sie unbehandelt zu verarbeiten.
    \item Sie will wissen, ob der Reissalat mit abgekochten Pfifferlingen besser schmeckt.
\end{itemize}

\subsection*{20. FV1.6: Ehh, roh, doch müssen se en bischen in de Pfanne tun.}

\textbf{Paraphrase:} Die Verkäuferin sagt, dass Pfifferlinge vor der Verarbeitung mit Hitze behandelt werden müssen.

\textbf{Intentionen:}
\begin{itemize}
    \item Sie will warnen (Gefahr).
    \item Sie will einen Zubereitungstipp geben.
\end{itemize}

\textbf{Regulative Regel:} Sind Gefahren mit der Ware verbunden, die dem Käufer nicht offensichtlich sind, muss der Verkäufer aufmerksam machen.

\subsection*{21. FK1.9: Tuh ich.}

\textbf{Paraphrase:} Die Käuferin sagt, dass sie die Pfifferlinge in die Pfanne tun wird.

\textbf{Intention:} Sie quittiert den Rat der Verkäuferin.

\subsection*{22. FV1.7: Klein bischen.}

\textbf{Paraphrase:} Die Verkäuferin rät, die Pfifferlinge ein wenig in der Pfanne zu erhitzen.

\textbf{Intention:} Sie bekräftigt ihren Rat.

\subsection*{24. FK1.10: Die kann ich aber, ehm, in en Reissalat tun.}

\textbf{Paraphrase:} Die Käuferin will erneut wissen, ob sie die Pfifferlinge nach der Behandlung in Reissalat tun kann.

\textbf{Intention:} Sie ist unsicher (vielleicht wegen Alters oder psychischer Instabilität).

\textbf{Regulative Regel:} Auf Fragen älterer oder psychisch instabiler Personen antwortet man geduldig.

\subsection*{25. FV1.8: Ja, datt is kein Problem, se müssen so nur...}

\textbf{Paraphrase:} Die Verkäuferin bestätigt erneut, dass es kein Problem ist.

\textbf{Intention:} Sie hat die Käuferin als älter oder psychisch instabil eingestuft.

\subsection*{26. FK1.11: bischen, ja}

\textbf{Paraphrase:} Die Käuferin sagt, dass sie die Pfifferlinge ein wenig in der Pfanne erwärmen wird.

\textbf{Intention:} Sie bestätigt, dass sie die Verkäuferin verstanden hat.

\subsection*{27. FV1.9: Bischen in eh, nitt wie de Champignons, die tuh ich ja auch roh erein.}

\textbf{Paraphrase:} Die Verkäuferin sagt, dass sie im Gegensatz zu Champignons die Pfifferlinge nur erwärmt in den Reissalat tun würde.

\textbf{Intention:} Sie ermahnt die Käuferin durch ein Gegenbeispiel.

\subsection*{28. FK1.11: eh ja.}

\textbf{Paraphrase:} Die Käuferin sagt, dass sie die Pfifferlinge erwärmen und dann erst in den Reissalat tun wird.

\textbf{Intentionen:}
\begin{itemize}
    \item Sie bestätigt den Rat.
    \item Sie merkt, dass sie wie ein Kind behandelt wird, und fordert ihre Zurechnungsfähigkeit ein.
\end{itemize}

\subsection*{29. FV1.10: Hundert ne?}

\textbf{Paraphrase:} Die Verkäuferin fragt, ob sie hundert Gramm Pfifferlinge abwiegen soll.

\textbf{Intentionen:}
\begin{itemize}
    \item Sie will endlich zum Verkaufsabschluss kommen.
    \item Sie quittiert die Einforderung von Zurechnungsfähigkeit mit einer sachlichen Frage.
\end{itemize}

\subsection*{30. FK1.12: Ja bitte. Watt krisch ich denn noch hier?}

\textbf{Paraphrase:} Die Käuferin bestätigt die Menge und fragt nach weiteren Angeboten.

\textbf{Intention:} Sie bestätigt und fragt aktiv nach weiteren Produkten.

\subsection*{31. FV1.11: Waldbeeren? Hab ich auch schonn.}

\textbf{Paraphrase:} Die Verkäuferin macht auf ihr Waldbeerangebot aufmerksam.

\textbf{Intention:} Sie macht ein weiteres Kaufangebot.

\subsection*{32. FK1.13: (unverständlich) Wie ist es denn mit Erdbeeren?}

\textbf{Paraphrase:} Die Käuferin fragt nach Erdbeeren.

\textbf{Intention:} Sie möchte keine Waldbeeren, sondern Erdbeeren.

\subsection*{34. FK1.14: Watt hann se denn sons noch?}

\textbf{Paraphrase:} Die Käuferin fragt nach weiteren Angeboten.

\textbf{Intention:} Sie hat ihren Wunsch nach Erdbeeren aufgegeben und fragt nach Alternativen.

\subsection*{35. FV1.12: Hann se denn keine Lust auf Himbeeren? Oder Johannisbeeren, hab ich auch schonn.}

\textbf{Paraphrase:} Die Verkäuferin bietet Himbeeren oder Johannisbeeren an.

\textbf{Intention:} Sie will die "Scharte" mit den Erdbeeren ausmerzen.

\subsection*{36. FK1.15: Ja. (Pause) Nehm werr beides eins.}

\textbf{Paraphrase:} Die Käuferin sagt, dass sie sowohl Himbeeren als auch Johannisbeeren nehmen will.

\textbf{Intention:} Sie will das Angebot annehmen und beide Sorten nehmen.

\subsection*{37. FV1.13: Johannisbeeren is a Pfund, die können se auch noch länger verwahren.}

\textbf{Paraphrase:} Die Verkäuferin weist auf die Pfundbindung und bessere Haltbarkeit hin.

\textbf{Intention:} Sie gleicht den Mangel der Pfundbindung mit dem Hinweis auf Haltbarkeit aus.

\subsection*{39. FK1.16: Dann habb ich, jlaub ich, alles fürr ze Hause.}

\textbf{Paraphrase:} Die Käuferin stellt fest, dass sie alles für zu Hause erworben hat.

\textbf{Intention:} Sie will signalisieren, dass sie ihre Kaufwünsche befriedigt hat.

\textbf{Regulative Regel:} Wer einen Kauf beenden will, muss das auch sagen.

\subsection*{40. FV1.14: Joh, bis Übermojen, nah.}

\textbf{Paraphrase:} Die Verkäuferin bestätigt das Ende und verweist auf den übermorgigen Tag.

\textbf{Intentionen:}
\begin{itemize}
    \item Sie will die Käuferin verabschieden.
    \item Sie will Kundenbindung erreichen.
\end{itemize}

\subsection*{41. FK1.17: neh. (Pause) Kuck mal, der junge Mann muß für Euch sorgen.}

\textbf{Paraphrase:} Die Käuferin macht auf eine männliche Person aufmerksam.

\textbf{Intention:} Sie will auf ein sie verblüffendes Ereignis hinweisen.

\subsection*{43. FK1.18: Ja}

\textbf{Paraphrase:} Die Käuferin bestätigt.

\textbf{Intention:} Sie will die Verkäuferin ermutigen, ihre Erklärung auszusprechen.

\subsection*{44. FV1.16: (unverständlich) damit uns ett Jehirrn nett ahfängt zu koche.}

\textbf{Paraphrase:} Die Verkäuferin erklärt, dass die Handlung des jungen Mannes verhindert, dass ihre Gehirne zu kochen anfangen.

\textbf{Intention:} Sie will die Handlung erklären.

\subsection*{45. FK1.19: So.}

\textbf{Paraphrase:} Die Käuferin sagt "So".

\textbf{Intention:} Sie will wieder zur Sache kommen.

\subsection*{46. FV1.17: Sechzig, vier Mark sechzig, acht Mark sechzig, zwölf Mark un fünfzig.}

\textbf{Paraphrase:} Die Verkäuferin rechnet den Preis aus.

\textbf{Intention:} Sie will zur Sache kommen und nennt den Endbetrag.

\subsection*{47. FK1.20: Du kriss die Tür nich zu.}

\textbf{Paraphrase:} Die Käuferin bringt ihr Erstaunen über die Summe zum Ausdruck.

\textbf{Intentionen:}
\begin{itemize}
    \item Sie will sagen, dass ihr der Preis zu hoch erscheint.
    \item Sie will sagen, dass sie erstaunt ist, wie viel sie gekauft hat.
\end{itemize}

\subsection*{48. FV1.18: Zwölf Mark un Fünfzisch. (Pause) Ich weiß, ich bin heut wieder unverschämt...}

\textbf{Paraphrase:} Die Verkäuferin wiederholt den Preis und bezeichnet sich als unverschämt.

\textbf{Intentionen:}
\begin{itemize}
    \item Sie bekräftigt ihre Forderung.
    \item Sie ironisiert den Einwand der Käuferin.
\end{itemize}

\subsection*{49. FK1.21: Ja.}

\textbf{Paraphrase:} Die Käuferin sagt "Ja".

\textbf{Intentionen:}
\begin{itemize}
    \item Sie bestätigt den Preis.
    \item Sie bestätigt, dass die Verkäuferin eine unverschämte Forderung stellt.
\end{itemize}

\subsection*{50. FV1.19: Aber, aber, aber, eine Mark (unverständlich) noch.}

\textbf{Paraphrase:} Die Verkäuferin beginnt mit einem Einwand, wechselt dann zu einer Restforderung.

\textbf{Intention:} Der Einwand ist überflüssig geworden, weil die Käuferin mit der Zahlung begonnen hat.

\subsection*{51. FK1.22: Hör ens}

\textbf{Paraphrase:} Die Käuferin fordert die Verkäuferin auf, ihr zuzuhören (wieder Du-Form).

\textbf{Intention:} Sie kündigt eine wichtige Mitteilung an.

\subsection*{52. FV1.20: Watt müssen se?}

\textbf{Paraphrase:} Die Verkäuferin fragt in der Sie-Form, worum es geht.

\textbf{Intention:} Sie will wissen, worum es geht, und unterstreicht persönliche Distanz.

\subsection*{53. FK1.23: Zur eh Barmer, aber ich komm dann, ich komm dann nachher, dann stell ich et unter.}

\textbf{Paraphrase:} Die Käuferin teilt mit, dass sie noch zur Krankenkasse muss.

\textbf{Intention:} Sie bittet um Hilfe.

\subsection*{54. FV1.21: neh, sons lassen se et hier. Dreizehn, fünfzhen, Zwanzig Mark.}

\textbf{Paraphrase:} Die Verkäuferin bietet an, die Ware zu verwahren, und gibt Wechselgeld zurück.

\textbf{Intention:} Man beendet den Kauf korrekt und freundlich.

\subsection*{55. FK1.24: Danke.}

\textbf{Paraphrase:} Die Käuferin dankt.

\subsection*{56. FV1.22: Bis Übermorgen.}

\textbf{Paraphrase:} Die Verkäuferin verabschiedet sich.

\subsection*{57. FK1.25: Danke schön.}

\textbf{Paraphrase:} Die Käuferin bedankt sich noch einmal.

\subsection*{58. FV1.23: Ja}

\textbf{Paraphrase:} Die Verkäuferin bekräftigt den Abschied und Dank.

\subsection*{59. *** Ende Text4 ***}

\textbf{Paraphrase:} Der Text ist beendet.

\textbf{Intention:} Das Ende des Textes wird markiert.

\newpage

\section{Tabellarische Übersicht aller Interakte}

\begin{longtable}{p{0.8cm}|p{3.5cm}|p{3.5cm}|p{3cm}|p{3cm}}
\hline
\textbf{Nr.} & \textbf{Interakt} & \textbf{Paraphrase} & \textbf{Intention(en)} & \textbf{Regeln} \\
\hline
1 & *** Anfang Text4 *** & Anfang eines Textes wird markiert & Leser soll wissen: neuer Text beginnt & Konstitutiv: Text braucht Anfang und Ende \\
2 & : Markt, 11.00 Uhr... & Historische/funktionale Einordnung & Kontext Markt-Institution & Regulativ: Nur zahlungsfähige Kunden erhalten Ware; Konstitutiv: Ohne Markt kein Güteraustausch \\
3 & (unverständlich) & Gespräch läuft schon & Zeigt: laufendes Gespräch wird verfolgt & – \\
4 & FK1.1: Hör ens, ich nehm ein paar Champignons mit. & Käuferin (Du) kündigt Mitnahme an & (1) einfach nehmen; (2) testen; (3) kaufen & – \\
5 & FV1.1: Eh, braune? & Frage nach braunen Champignons & Bedarf klären & Regulativ: Bei Unklarheit nachfragen; Konstitutiv: Verkaufsgespräch erfordert Bedarfsklärung \\
6 & FK1.2: Ne helle. & Spezifikation: helle Champignons & Bedarf spezifizieren & – \\
7 & FV1.2: Helle. & Bestätigung: helle Champignons & (1) geben; (2) probieren lassen; (3) verkaufen & – \\
8 & FK1.3: mhmh. & Bestätigung & Veranlasst Verkäuferin zu handeln & – \\
9 & (unverständlich) & – & – & – \\
10 & FK1.4: Meinen se nich. & Wechsel in Sie-Form: Frage nach Meinung & (1) nehmen; (2) testen; (3) kaufen & – \\
11 & FV1.3: Ja is ejal, se sinn beide frisch. & Beide Sorten frisch & Entscheidung erleichtern / Bedenken zerstreuen & – \\
12 & FK1.5: Oder, wie is et denn mit, mit, eh & Unsicherheit, neuer Wunsch & – & Regulativ: Kaufwunsch darf in Grenzen geändert werden; Konstitutiv: Flexibilität ist Erfolgsbedingung \\
\hline
\multicolumn{5}{l}{\textit{Ab hier werden Nicht-Kauf-Lesarten aufgegeben.}} \\
\hline
13 & FV1.4: Die können se länger liejen lassen. & Grund: helle Champignons lagern länger & (1) echter Rat; (2) abgesicherte Täuschung & Regulativ: Auf Unsicherheiten eingehen; Konstitutiv: Verkaufsgespräch = Beratungsgespräch \\
14 & FK1.6: Neh, aber Pfifferlinge. & Frage nach Pfifferlingen & Kaufwunsch Pfifferlinge & Regulativ: Kaufwunsch signalisieren; Konstitutiv: Ohne Kaufwunsch kein Kauf \\
15 & FV1.5: Ah, die sinn super. & Begeisterung über Pfifferlinge & Produkt positiv hervorheben & – \\
17 & FK1.7: Kann ich die denn in Reissalat tun? & Frage zur Verwendung & Wissen, ob es geht & Regulativ: Fragen erlaubt; Konstitutiv: Fragen sind konstitutiv für Verkaufsgespräche \\
19 & FK1.8: Brauch ich nich abzukochen oder was? & Frage zur Zubereitung & (1) Gefahr; (2) Geschmack & – \\
20 & FV1.6: Ehh, roh, doch müssen se en bischen in de Pfanne tun. & Hitze-Behandlung nötig & (1) warnen; (2) Tipp geben & Regulativ: Bei Gefahr aufklären \\
21 & FK1.9: Tuh ich. & Bestätigung & Rat quittieren & – \\
22 & FV1.7: Klein bischen. & Präzisierung: wenig erhitzen & Rat bekräftigen & – \\
24 & FK1.10: Die kann ich aber, ehm, in en Reissalat tun. & Erneute Frage & Unsicherheit (Alter/Psyche) & Regulativ: Geduld bei älteren/psychisch instabilen Kunden \\
25 & FV1.8: Ja, datt is kein Problem, se müssen so nur... & Erneute Bestätigung & Einstufung als instabil – geduldige Behandlung & – \\
26 & FK1.11: bischen, ja & Bestätigung & Verstanden haben & – \\
27 & FV1.9: Bischen in eh, nitt wie de Champignons... & Gegenbeispiel Champignons & Ermahnung & Regulativ: psychisch Instabile wie beschränkt Geschäftsfähige behandeln \\
28 & FK1.11: eh ja. & Bestätigung + leichter Widerstand & (1) Rat bestätigen; (2) Zurechnungsfähigkeit einfordern & Konstitutiv: Erwachsene fordern gegenseitig Zurechnungsfähigkeit ein \\
29 & FV1.10: Hundert ne? & Frage nach Menge & (1) Abschluss herbeiführen; (2) sachlich bleiben & – \\
30 & FK1.12: Ja bitte. Watt krisch ich denn noch hier? & Bestätigung + Frage nach weiteren Produkten & Aktive Suche nach weiteren Kaufmöglichkeiten & – \\
31 & FV1.11: Waldbeeren? Hab ich auch schonn. & Angebot Waldbeeren & Weiteres Kaufangebot & Konstitutiv: Käufer kaufen gern bei guten Verkäufern \\
32 & FK1.13: ... Wie ist es denn mit Erdbeeren? & Frage nach Erdbeeren & Keine Waldbeeren, sondern Erdbeeren & – \\
34 & FK1.14: Watt hann se denn sons noch? & Frage nach Alternativen & Erdbeer-Wunsch aufgegeben – Alternativsuche & – \\
35 & FV1.12: Hann se denn keine Lust auf Himbeeren? Oder Johannisbeeren... & Angebot Himbeeren/Johannisbeeren & "Scharte" der Erdbeeren ausmerzen & Regulativ: "Bluffen ist alles"; Konstitutiv: "Wer zögert, wird ausgetrixt" \\
36 & FK1.15: Ja. (Pause) Nehm werr beides eins. & Nimmt beide Sorten & Angebot annehmen & – \\
37 & FV1.13: Johannisbeeren is a Pfund... & Hinweis auf Pfundbindung und Haltbarkeit & Mangel ausgleichen & – \\
39 & FK1.16: Dann habb ich... alles fürr ze Hause. & Kauf abgeschlossen & Beenden signalisieren & Regulativ: Beenden muss angekündigt werden; Konstitutiv: Kaufhandlung hat ein Ende \\
40 & FV1.14: Joh, bis Übermojen, nah. & Bestätigung, Verweis auf übermorgen & (1) verabschieden; (2) Kundenbindung & Regulativ: Freundlichkeit; Konstitutiv: Verkäufer sind freundlich \\
41 & FK1.17: neh. (Pause) Kuck mal, der junge Mann muß für Euch sorgen. & Hinweis auf jungen Mann & Unerwartetes Ereignis teilen & Regulativ: Unerwartetes teilen; Konstitutiv: Geteilte Erfahrung erhöht Lebensgenuss \\
43 & FK1.18: Ja & Bestätigung & Ermutigung zur Erklärung & Regulativ: Im Gespräch bleiben; Konstitutiv: Nur wer im Gespräch bleibt, ist dabei \\
44 & FV1.16: ...damit uns ett Jehirrn nett ahfängt zu koche. & Erklärung der Handlung & Handlung erklären & Regulativ: Handlungen müssen erklärbar sein; Konstitutiv: Unerklärliche Handlungen verursachen Unsicherheit \\
45 & FK1.19: So. & "So" & Zur Sache kommen & Regulativ: Verkaufsgespräche müssen enden; Konstitutiv: Endlose Gespräche zerstören Fortgang \\
46 & FV1.17: Sechzig, vier Mark sechzig... & Preisberechnung & Endbetrag nennen & Regulativ: Preis in Landeswährung fordern; Konstitutiv: Nur mit Zahlung erfolgreicher Abschluss \\
47 & FK1.20: Du kriss die Tür nich zu. & Erstaunen über Summe & (1) Preis zu hoch; (2) Erstaunen über Menge & Regulativ: Preis drücken; Konstitutiv: Marktpreise entstehen durch Aushandlung \\
48 & FV1.18: Zwölf Mark... Ich weiß, ich bin heut wieder unverschämt... & Preis wiederholen, selbstironisch & (1) Forderung bekräftigen; (2) ironisieren & – \\
49 & FK1.21: Ja. & Bestätigung & (1) Preis bestätigen; (2) Unverschämtheit bestätigen & Konstitutiv: Über Preise muss man sich einigen \\
50 & FV1.19: Aber, aber, aber, eine Mark... noch. & Einwand, dann Restforderung & Zahlung hat begonnen – Restforderung & Regulativ: Preise vollständig bezahlen; Konstitutiv: Nur vollständige Bezahlung = Kaufvertrag \\
51 & FK1.22: Hör ens & Aufforderung zuzuhören (Du) & Wichtige Mitteilung ankündigen & – \\
52 & FV1.20: Watt müssen se? & Frage (Sie) & Wissen, worum es geht – Distanz wahren & – \\
53 & FK1.23: Zur eh Barmer... & Mitteilung: Krankenkassenbesuch & Um Hilfe bitten & – \\
54 & FV1.21: neh, sons lassen se et hier. Dreizehn, fünfzhen, Zwanzig Mark. & Verwahren anbieten + Wechselgeld & Korrekter, freundlicher Abschluss & – \\
55 & FK1.24: Danke. & Dank & Dank für Hilfe/Abschluss & – \\
56 & FV1.22: Bis Übermorgen. & Verabschiedung & – & – \\
57 & FK1.25: Danke schön. & Erneuter Dank & – & – \\
58 & FV1.23: Ja & Bekräftigung & – & – \\
59 & *** Ende Text4 *** & Textende markieren & Ende markieren & Regulativ: Textende markieren; Konstitutiv: Texte haben endliche Länge \\
\hline
\end{longtable}

\newpage

\section{Visualisierung: Der sich verengende Lesarten-Trichter}

\begin{center}
\begin{tikzpicture}[
    node distance=1cm and 1.5cm,
    every node/.style={align=center, font=\small},
    level/.style={thick, ->},
    trim/.style={draw, rectangle, rounded corners, minimum width=8cm, minimum height=0.8cm},
    entscheidung/.style={draw, diamond, aspect=2, inner sep=2pt, font=\tiny}
]

% Hintergrund-Trichter (visuelle Andeutung)
\fill[gray!10] (-5,7) -- (5,7) -- (2.5,-2) -- (-2.5,-2) -- cycle;
\draw[gray!30, dashed] (-5,7) -- (-2.5,-2);
\draw[gray!30, dashed] (5,7) -- (2.5,-2);

% Ebene 1: Beginn (maximale Lesartenvielfalt)
\node[trim, fill=blue!10] (e1) at (0,6.5) {Interakt 4: "Hör ens, ich nehm ein paar Champignons mit"\\\textbf{3 Lesarten}: (1) einfach nehmen / (2) probieren / (3) kaufen};

% Ebene 2: Erste Einschränkung durch Fortsetzung
\node[trim, fill=yellow!15] (e2) at (0,4.5) {Interakt 5-8: Bedarfsklärung (FV1.1: "Eh, braune?" / FK1.2: "Ne helle" / FV1.2: "Helle" / FK1.3: "mhmh")\\\textbf{Alle drei Lesarten noch möglich, aber Kontext verengt}};

% Entscheidungspfeil
\draw[level] (e1) -- (e2);

% Ebene 3: Kritische Weichenstellung
\node[trim, fill=orange!20] (e3) at (0,2.5) {Interakt 10-12: Wechsel in Sie-Form + Unsicherheit\\"Meinen se nich" / "Oder, wie is et denn mit..."\\\textbf{Nur noch Kauf-Lesarten plausibel}};

\draw[level] (e2) -- (e3);

% Ebene 4: Finale Einengung
\node[trim, fill=red!20] (e4) at (0,0.5) {Interakt 14: "Neh, aber Pfifferlinge"\\\textbf{Nur noch: konkreter Kaufwunsch}};

\draw[level] (e3) -- (e4);

% Ergebnis
\node[draw, rectangle, rounded corners, thick, fill=green!20, minimum width=6cm, minimum height=0.8cm] (ergebnis) at (0,-1.5) {Vorläufig bewährte Struktur:\\Kaufwunsch (Typ 3: VKG mit B-Schleife)};

\draw[level] (e4) -- (ergebnis);

% Seitenmarkierungen: Verworfenes
\node[font=\tiny, text=gray] (verw1) at (6,6.5) {↗ später verworfen};
\node[font=\tiny, text=gray] (verw2) at (6,4.5) {↗ später verworfen};
\node[font=\tiny, text=gray] (verw3) at (6,2.5) {✗ Nicht-Kauf-Lesarten\\(z.B. einfach nehmen,\\probieren)};
\node[font=\tiny, text=gray] (verw4) at (6,0.5) {✗ Unbestimmter Kaufwunsch};

% Legende unten
\draw[gray, dashed] (-5,-2.5) -- (5,-2.5);
\node[font=\small, below] at (0,-3.2) {
\begin{tabular}{ll}
\textcolor{blue!80}{●} & Maximale Lesartenvielfalt (zu Beginn) \\
\textcolor{orange!80}{●} & Kritische Weichenstellung (Lesarten werden ausgeschlossen) \\
\textcolor{red!80}{●} & Finale Einengung (nur noch eine Lesart plausibel) \\
\textcolor{green!80}{●} & Vorläufig bewährtes Ergebnis \\
\end{tabular}
};
\end{tikzpicture}
\end{center}

\end{document}